\begin{abstract}
Commuters driving into Hsinchu Science Park on National Freeway~1 decide each morning whether
to leave now or later. We forecast the car travel time of the 4.8\,km Zhubei--Hsinchu segment
60 minutes ahead from 20 months of 5-minute toll-gantry data, using only information available
at prediction time. A Ridge regression adds upstream travel time, calendar flags and
workday$\times$\allowbreak time-slot interactions to the current travel time. On a held-out
summer its peak-hour MAE is \res{test.ridgeC.peak}\,min (95\% CI
\res{test.peakLo}--\res{test.peakHi}), within our \res{cfg.threshold}-minute decision threshold
and \res{boot.persistence.peak.diff}\,min below persistence; the slot interactions are the most
valuable feature group. Gradient boosting is better overall (test MAE \res{test.lgbm.mae} vs.\
\res{test.ridgeC.mae}\,min) but worse in the peak. In routine peaks the model is only
\res{boot.profile.peak.diff}\,min better than a historical per-slot average (CI includes 0); its
value lies in rain and congestion onset.
\end{abstract}
